% Generated draft: requires mathematical and source review.
\documentclass[11pt]{article}
\usepackage[margin=25mm]{geometry}
\usepackage{fontspec}
\setmainfont{Latin Modern Roman}
\usepackage{amsmath,amssymb,mathtools,unicode-math}
\setmathfont{Latin Modern Math}
\usepackage{enumitem,needspace,xurl,hyperref}
\setlength{\parindent}{0pt}
\setlength{\parskip}{5pt}
\setlength{\emergencystretch}{3em}
\newcommand{\N}{\mathbb{N}}
\newcommand{\Z}{\mathbb{Z}}
\newcommand{\Q}{\mathbb{Q}}
\newcommand{\R}{\mathbb{R}}
\newcommand{\C}{\mathbb{C}}
\newcommand{\F}{\mathbb{F}}
\newcommand{\A}{\mathbb{A}}
\newcommand{\PP}{\mathbb{P}}
\newcommand{\E}{\mathbb{E}}
\newcommand{\eps}{\varepsilon}
\newcommand{\dd}{\,\mathrm{d}}
\begin{document}
\section*{UnsolvedMath ID: 800789}
\textbf{Source problem code:} catalogue entry
\subsection{Definitions and mathematical statement}
Let \(n\geq 1\), let \(O\) denote the origin of \(\C^n\), and let \(\operatorname{End}(\C^n,O)\) denote the holomorphic map germs \(f:(\C^n,O)\to(\C^n,O)\) fixing \(O\). A germ is superattracting if \(Df(O)=0\). Let \(f\) be such a germ with a first nonzero homogeneous term, and write
\[
f(z)=H(z)+\sum_{j>k}H_j(z),\qquad k\geq 2,
\]
where each \(H_j\) is a homogeneous polynomial map of degree \(j\), and \(H\not\equiv 0\) is homogeneous of degree \(k\). Assume that \(H\) is degenerate in the precise sense
\[
\exists v\in\C^n\setminus\{O\}\quad H(v)=O.
\]
The critical locus is \(\operatorname{Crit}(f)=\{z:\det Df(z)=0\}\), interpreted locally. Write \(\mathcal P_f\) for the postcritical set intended by the source. The usual forward-image expression is
\[
\bigcup_{m\geq 1} f^{\circ m}\bigl(\operatorname{Crit}(f)\bigr),
\]
with images taken where the iterates are defined. The supplied record does not specify whether this expression, its closure, or another local interpretation is intended; thus the convention for \(\mathcal P_f\) remains unresolved. An analytic cone at \(O\) means a complex analytic set germ invariant under complex scalar multiplication, understood as a germ.

\textbf{Question (Abate, Problem OP7).} For every germ \(f\) satisfying the preceding hypotheses, if its postcritical set \(\mathcal P_f\) is an analytic cone, must there exist a biholomorphic germ \(\Phi:(\C^n,O)\to(\C^n,O)\) such that
\[
\Phi\circ f\circ\Phi^{-1}=H?
\]
This formulation retains the source's unresolved postcritical-set convention. No finiteness, dominance, or isolated-zero hypothesis on \(f\) or \(H\) is supplied.
\subsection{Short English statement}
Abate asks whether a superattracting holomorphic germ whose first nonzero homogeneous term vanishes at a nonzero vector must be locally holomorphically conjugate to that term whenever its postcritical set is an analytic cone. The precise local convention for the postcritical set requires clarification.
\subsection{Sources}
\begin{itemize}
\item Local UnsolvedMath registry, permanent ID \(800789\), source problem code \texttt{catalogue entry}, title ``Local Discrete Holomorphic Dynamics --- Open Problem OP7.'' The supplied record attributes the problem to Abate, with proposed year 2011, and labels it \texttt{open}; that label is supplied registry information, not an independently verified status assessment.
\item Abate, \emph{Open Problems in Local Discrete Holomorphic Dynamics} (2011), Problem OP7. Supplied source URL: \url{https://arxiv.org/abs/1107.4876}. The background records access on 2026-10-16 and extraction by \texttt{arxiv\_latex}; these are supplied provenance claims, not access or verification performed for this draft.
\item Supplied upstream record ID \(800789\), revision \texttt{d371ee917ac9647fbb06cd6fc529deec6f3dda82}; registry creation and update dates are both 2026-10-16.
\end{itemize}
\end{document}
